% !PRINTED 165
\weberpage{165}
In der Tabelle V sind die vier Formeln, die sich auf diese Weise ergeben, zusammengestellt. Subtrahiert man die Formel (3) dieser Tabelle von den drei übrigen, so ergeben sich Entwickelungen für die logarithmischen Ableitungen der drei elliptischen Grundfunktionen, von denen wir die eine anführen wollen:
\weq{4}{\frac{d\log\operatorname{sn}2Ku}{\pi\,du} = \frac{d\log\sin\pi u}{\pi\,du} + 4\sin 2\pi u\sum_{h=1}^{\infty}\frac{(-1)^h q^h}{1-2q^h\cos 2\pi u+q^{2h}}}
\[
= \frac{d\log\sin\pi u}{\pi\,du} - 4\sum_{h=1}^{\infty}\frac{q^h\sin 2h\pi u}{1+q^h}.
\]

Entwickelungen für die Transzendenten \gesperrt{dritter Gattung} ergeben sich, wenn man die Formeln der fünften Tabelle zwischen den Grenzen $u+v$, $u-v$ integriert. Man erhält so z.\,B. aus der Formel (3) dieser Tabelle:
\weq{5}{\frac{1}{2}\log\frac{\theta[2K(u-v)]}{\theta[2K(u+v)]} = -4\sum\frac{q^{\frac{m}{2}}}{m(1-q^m)}\sin m\pi u\sin m\pi v.}

\begin{center}
\rule{2cm}{0.4pt}
\end{center}

Setzt man in den Formeln der Tabelle IV und V $u=0$, so ergeben sich die folgenden Entwickelungen:
\[
\begin{aligned}
\frac{2K}{\pi} &= 1+4\sum\frac{q^{\frac{m}{2}}}{1+q^m} \phantom{,} = 1+4\sum\frac{(-1)^{\frac{n-1}{2}}q^n}{1-q^n}, \\
\frac{2\varkappa K}{\pi} &= \phantom{1+{}} 4\sum\frac{q^{\frac{n}{2}}}{1+q^n} \phantom{,} = \phantom{1+{}} 4\sum\frac{(-1)^{\frac{n-1}{2}}q^{\frac{n}{2}}}{1-q^n}, \\
\frac{2\varkappa'K}{\pi} &= 1+4\sum\frac{(-1)^{\frac{m}{2}}q^{\frac{m}{2}}}{1+q^m} = 1-4\sum\frac{(-1)^{\frac{n-1}{2}}q^n}{1+q^n}.
\end{aligned}
\]

Endlich erhält man eine Entwickelung für das vollständige Integral zweiter Gattung $E$ aus der Formel
\[
Z(v) = E(v) - \frac{E}{K}v.
\]
Setzt man hierin nach der Differentiation $v=0$, so folgt:
\[
Z'(0) = \frac{K-E}{K},
\]
und mithin nach der Tabelle V, Formel (3):
\[
K-E = \frac{2\pi^2}{K}\sum\frac{q^n}{(1-q^n)^2} = \frac{\pi^2}{K}\sum\frac{mq^{\frac{m}{2}}}{1-q^m}.
\]

\begin{center}
\rule{2.8cm}{0.4pt}
\end{center}
